\section{Fundamentación Teórica}

A. Programación Orientada a Objetos (POO)

La Programación Orientada a Objetos organiza el software en torno a clases y objetos, favoreciendo principios como la abstracción, el encapsulamiento, la herencia y el polimorfismo. Según Introducción a la Programación Orientada a Objetos \cite{poo}, este paradigma promueve la reutilización del código y facilita la modularidad de los proyectos. En comparación con la programación estructurada, la POO ofrece mayor claridad y un modelado más cercano a la lógica del mundo real. La abstracción permite modelar conceptos del entorno dentro del software, mientras que el encapsulamiento protege y mantiene la coherencia de los datos internos. La herencia fomenta la reutilización del código al establecer jerarquías entre clases, y el polimorfismo facilita la extensión de funcionalidades sin alterar estructuras ya establecidas. Además, se resalta el valor de la POO en la enseñanza de la programación, ya que brinda a los estudiantes una forma intuitiva de abordar problemas complejos y mantiene su vigencia en lenguajes como Java, C\#, Python o C++ \cite{poo}.

B. Patrón Modelo–Vista–Controlador (MVC)

El patrón Modelo–Vista–Controlador (MVC) divide una aplicación en tres componentes principales: el Modelo (datos y lógica de negocio), la Vista (interfaz de usuario) y el Controlador (intermediario). De acuerdo con Automatización de la codificación del patrón MVC en proyectos orientados a la Web \cite{mvc}, esta separación mejora la productividad, la escalabilidad y la consistencia del software, además de facilitar pruebas unitarias y mantenimiento. La propuesta técnica de dicho trabajo presenta una herramienta en Java (GAPRIP MVC) que automatiza la codificación a partir de esquemas de bases de datos en MySQL. El sistema genera clases DTO y DAO, junto con capas de negocio y presentación, lo que reduce tiempos de desarrollo en un 60–70 \%. Esto confirma que el patrón MVC, además de organizar el código, permite integrar procesos de automatización que potencian la eficiencia en proyectos web \cite{mvc}.

C. Node.js y Spring Boot

Node.js es un entorno ligero y eficiente, ideal para aplicaciones en tiempo real y microservicios, mientras que Spring Boot se caracteriza por su robustez y seguridad en entornos empresariales basados en Java. El estudio Node.js vs Spring Boot en el desarrollo backend de aplicaciones web \cite{nodejs_springboot} concluye que la elección depende de las necesidades del proyecto: Node.js destaca por su flexibilidad y escalabilidad, mientras que Spring Boot ofrece solidez y estabilidad. En particular, Spring Boot simplifica la construcción de microservicios al proveer autoconfiguración y servidores embebidos, lo que disminuye la complejidad del despliegue. Además, su madurez le permite integrarse con librerías como Spring Security y bases de datos SQL transaccionales. Por otro lado, Node.js, junto con Express.js, aprovecha el modelo asíncrono de JavaScript para ofrecer rapidez y facilidad de integración con bases NoSQL como MongoDB, aunque con ciertas limitaciones en materia de seguridad y dependencia de paquetes externos \cite{nodejs_springboot}.

D. Metodología Scrum

Scrum se basa en iteraciones cortas llamadas Sprints y define roles esenciales como el Scrum Master y el Product Owner. Según Metodología Scrum: Desarrollo detallado de la fase de aprobación de un proyecto informático mediante el uso de metodologías ágiles \cite{scrum}, este marco de trabajo mejora la organización del equipo, la comunicación con el cliente y la entrega continua de valor. El documento explica cómo los elementos clave —Product Backlog, Sprint Backlog e Incrementos de producto— permiten priorizar requerimientos y entregar software funcional en cada ciclo. Además, su carácter iterativo y adaptable refuerza la capacidad de respuesta frente a cambios del entorno, reduciendo riesgos y aumentando la satisfacción del cliente \cite{scrum}.

E. Bases de datos

Las bases de datos son el núcleo en la gestión de información. El libro Fundamentos sobre la gestión de base de datos \cite{databases} explica que los modelos relacionales, apoyados en SQL, aseguran consistencia e integridad mediante procesos de normalización y control de dependencias. Por su parte, los modelos NoSQL ofrecen flexibilidad para manejar grandes volúmenes de datos, lo que resulta útil en aplicaciones modernas distribuidas y de rápido crecimiento \cite{databases}. En general, la gestión de bases de datos permite construir sistemas robustos y confiables. Operaciones como consultas, inserciones y actualizaciones garantizan precisión y disponibilidad, mientras que los mecanismos de control de acceso y llaves primarias y foráneas protegen frente a inconsistencias. Estos principios consolidan a las bases de datos como un componente esencial en cualquier sistema de información.
